\section{n1. Whole410 canonical envelope filtration independently reconstructed on every excellent reduced germ; general continuity remains unproved.}
\textit{Supports: G331.}

Relative only to the ENTIRE qualified G331, let R be excellent local, replace it by Rred, and fix a proper primary ideal a. Its positive-dimensional section X=V\_a(R) includes EVERY support; dimension zero has the unique empty envelope. Write C for the continuous independently geometric class and D for its actual decreasing-net completion. G331 identifies these with the corresponding ideal classes, after defining positivity by actual negative effective real Cartier presentations and every contracted-curve degree test on directed actual primary blowups. This comparison, rational profile density, finite maxima, and full-fibre positive descent are used with every source/scope qualification. No envelope theorem is imported.

For finite continuous signed f define A\_0=R and A\_m=\{g: v(g)+m f(v)\textgreater{}=0 for ALL v in X\}. If c=max(0,-min f), then a\textasciicircum{}N is contained in A\_m for every positive N\textgreater{}=mc. Ultrametric addition and nonnegative ring values make A\_m an ideal. If f is negative somewhere no unit satisfies the defining inequality, so every positive-degree A\_m is proper primary. Otherwise A\_m=R. Multiplication gives A\_m A\_n subset A\_(m+n). Every A\_m is integrally closed: in a monic ideal-integrality equation, v(g)\textless{}v(A\_m) would make the leading term uniquely least of finite value. Infinite v(g) causes no exception. The finite continuous profile h\_m=-v(A\_m)/m lies below f and is geometrically positive: -E\_(A\_m)/m is actual negative effective Cartier and relatively ample on its actual primary blowup. Unit profiles are zero.

Define the RAW function p(f)=sup\{u in D:u\textless{}=f\}, without asserting positivity of this supremum. It equals the supremum over C and sup\_m h\_m. Indeed a rational competitor -v(I)/r\textless{}=f has I subset A\_r. A C competitor is uniformly approximated by rational profiles; shift an approximant down by a small rational a-constant, realized by ideal multiplication and a common denominator, to obtain rational competitors arbitrarily close from below. For u=inf\_i g\_i\textless{}=f in an ACTUAL decreasing C-net, the open sets \{g\_i\textless{}f+epsilon\} cover compact X and are directed by inclusion. One index works everywhere. Then g\_i-epsilon is a C competitor and u\textless{}=sup\_m h\_m+epsilon. Let epsilon decrease to zero. This handles bottom and partial poles without continuity of u.

Set M\_1=||v(A\_1)|| and H\_N=max\_(1\textless{}=r\textless{}=N)h\_r in C. For m=kr+s, 0\textless{}=s\textless{}r, grading gives h\_m\textgreater{}=(kr/m)h\_r-sM\_1/m\textgreater{}=h\_r-rM\_1/m; hence H\_N-NM\_1/m\textless{}=h\_m\textless{}=p. Therefore the FULL h\_m sequence converges pointwise to the finite nonpositive lower-semicontinuous p. Moreover p is continuous iff H\_N converges uniformly iff h\_m converges uniformly. If p is continuous, the increasing strict-error open sets for H\_N cover compact X, proving uniform convergence; the residue bound transfers it to h\_m. Conversely a uniform limit lies in closed C and equals p. In that case p is the greatest D minorant, with actual geometric positivity. Pointwise convergence alone supplies no regularity.

The raw map is sup-norm 1-Lipschitz. If f\textless{}=g+epsilon, subtract the fixed negative primary constant epsilon from a D competitor below f to obtain a D competitor below g. Supremizing and reversing f,g proves the bound. Thus uniform density of a specified class of obstacles with C envelopes would prove CE for its uniform closure; no general such finite-generation density is asserted.

For this FIXED obstacle, its canonical algebra A=direct-sum A\_m t\textasciicircum{}m is finitely generated iff p is a rational unit-or-primary profile iff a finite maximum of canonical profiles attains p. To prove the forward direction, take homogeneous generators g\_i t\textasciicircum{}(d\_i), let L be divisible by all d\_i and D\_0=max d\_i, and put J=(g\_i\textasciicircum{}(L/d\_i)). At each support some generator has finite value because A\_1 contains a primary power. Degree-m expressions have value at least m min\_i v(g\_i)/d\_i, while pure powers belong to A\_L. Consequently p=h\_L=-v(J)/L. A generator attaining the minimum also shows p=max\_(r\textless{}=D\_0)h\_r. The radical of J contains A\_1, so J is primary in the proper case. The residue bound yields 0\textless{}=p-h\_m\textless{}=D\_0 M\_1/m uniformly on ALL supports.

Conversely suppose p=-v(J)/q with J primary and q positive integral. Since h\_m\textless{}=p\textless{}=f, every g in A\_m satisfies v(g)\textgreater{}=(m/q)v(J). The reverse implication follows from p\textless{}=f. Thus A\_m=\{g:g\textasciicircum{}q in integral-closure(J\textasciicircum{}m)\}. The equality uses the centred rank-one DVR criterion on every minimal-component domain quotient and the all-minimal-component criterion, precisely SH6.8.3--4 and 1.1.5. Every minimal component has positive dimension: otherwise the maximal ideal is minimal and the local reduction has dimension zero. Each centred DVR normalizes into X. Monic equations prove the reverse inequalities on every other support too. These complete immediate source units were freshly read, with formula bars checked in the actual PDF pixels; deeper valuation-domination foundations are named conventional background.

Let B=R[Jt\textasciicircum{}q]. A primary ideal contains a regular element by finite prime avoidance, so B is reduced excellent and has total quotient product K\_j(t\textasciicircum{}q). Its finite extension product K\_j(t) may be inseparable. Excellent/Nagata finite normalization makes the integral closure finite over B. An equation for g\textasciicircum{}q over J\textasciicircum{}m, multiplied by t\textasciicircum{}(mq) powers, makes g\textasciicircum{}q t\textasciicircum{}(mq) integral over B; the qth-power equation then makes gt\textasciicircum{}m integral. Hence A is a B-submodule of that finite normalization and is finite by Noetherianity. The unit case gives R[t]. A finite maximum of rational profiles is rational via sums of powered ideals. No arbitrary section algebra is assumed finitely generated.

For an integral signed exceptional Cartier obstacle f\_D on an actual primary model over normal R, let Y be its finite normalization. The canonical rational trivialization gives A\_m=H\textasciicircum{}0(Y,O(mD)) intersect R. All-boundary tests at exceptional height-one orders imply the section inequalities there; at other height-one primes g is regular and D vanishes. Normal height-one intersection gives section membership. Conversely pull the section to the valuation ring of Frac(R/p) for each support p. Properness gives its centre, and the generic point is outside the exceptional fibre. The invertible lattice has nonzero generic trivialization; if b generates it, v(g)\textgreater{}=v(b)=-m f\_D(v). This is the defining inequality on every support, with infinite g values automatic. No original fraction field is mapped through p. This computation supplies no general finite-generation theorem.

For another primary b put d(v)=v(b)\textgreater{}0 and T(v)=v/d(v), transporting f by f\_b(Tv)=f\_a(v)/d(v). The defining inequalities divide by d. ALL A\_m and their algebra are literally unchanged; h\_m and p divide by d, with the exact three-gauge cocycle. Uniform constants divide by min d\textgreater{}0. Finite normalization S gives A\_m\textasciicircum{}R=A\_m\textasciicircum{}S intersect R by surjective restriction. If a source envelope exists and is constant on EVERY fibre, G331 descends it positively; it dominates all pulled-back downstairs competitors, so is the downstairs envelope. Obstacle compatibility alone is insufficient, including on infinite-conductor quotient fibres. A decreasing family of continuous obstacles with C envelopes H\_i has the actual D envelope inf\_i H\_i of its infimum: it is a minorant and every competitor lies below each H\_i.

On every excellent reduced curve, normalization centres are DVRs and a regular conductor element has finite branch order. Thus boundary fibres are singletons. Actual real Cartier slopes on each DVR have no contracted-curve conditions and give every finite nonpositive coordinate. Semilocal gluing and G331 positive descent give C equal to the finite nonpositive orthant; decreasing coordinate truncations give the extended orthant D. Consequently the geometric envelope is min(f,0), matching the ideal envelope. This does not identify the valuative class with literal multibranch classical/CLD logarithmic radial extensions: the mixed-zero obstruction remains. Nilpotents are invisible and dimension-zero reduction has empty boundary. On a DVR the irrational obstacle -alpha has continuous real Cartier envelope, but A\_m=(x\textasciicircum{}ceil(m alpha)) is not finitely generated: finite homogeneous generators would force a rational minimum order/degree ratio equal to alpha. Thus non-finite-generation is not a CE counterexample.

This independently reconstructs the whole410 canonical argument, without using whole410 as a correctness premise. Artifacts: n1.txt, distinct TeX exposition n1.tex, P12.3-algebraic-descent-envelopes.tex, proof-companion.md and source-application-reading-ledger.md. Rerun own programs: python audit\_data.py; python assemble\_delivery.py; python check\_delivery.py. Mathematical claims are draft proofs, not consequences of execution.

\subsection*{Expanded canonical formulae and model qualification}
The residue degree estimate is valid for every $m,r\ge1$, including $m<r$:
\[
 A_r^kA_1^s\subset A_m,\quad m=kr+s,\quad0\le s<r,
 \qquad H_N-\frac{NM_1}{m}\le h_m\le p(f).
\]
Thus the compactness argument uses $H_N$, not a presumed monotone full
sequence.  Finite generation refers to this obstacle's algebra:
\[
 \bigoplus_{m\ge0}A_mt^m\text{ finitely generated}
 \ \Longleftrightarrow\ p(f)=-\frac1qv(J)
 \ \Longleftrightarrow\ p(f)=\max_{r\le N}h_r.
\]
In the normal section computation, height-one tests are made on the
finite normalization of the actual primary model, whereas geometric degree
positivity remains the actual primary-model test from G331.  Neither role
identifies arbitrary normal Weil coefficients with actual Cartier positivity.

